\chapter{How Can We Know an Unrecorded Past?}
\section{Half a Copper Coin}
First, pick something up: half a copper coin.

It lies at the very back of your grandfather's drawer, crowded among old keys, dead watches, and reading glasses with a broken arm. The coin is round, with a square hole in the middle. Its edge is chipped, and only a little more than half remains. There should have been four characters on its face, an era name followed by tongbao, meaning circulating treasure. Now the surface is so worn that you can barely recognize half of the character tong. You weigh it between your fingers. It is light. Green copper corrosion rubs onto your fingertips, carrying an earthy metallic smell.

You run to ask your grandfather, ``When is this from?'' He says he found it in a field as a child. Which dynasty it belongs to, he cannot say. ``An old thing, anyway.''

Notice the strange thing happening here. This coin cannot speak. It has no instruction manual, and no camera recorded who spent it or what it bought. Yet you, I, anyone who sees it, will agree without prompting that it proves \textbf{some things really existed in the past}. Someone cast it, someone used it, someone lost it. It is a small piece of evidence left by a history without recordings.

Now magnify the question ten thousand times. Not just this coin: the feelings of Shang people during divination, prices on the streets of ancient Rome, a day in the life of a Tang farmer's daughter, the sound of camel bells in a fourteenth-century West African caravan. None has a video, an audio recording, or a livestream replay. Where, then, do those clearly described ``pasts'' in textbooks come from? What entitles us to say we ``know''?

This chapter answers that question. The answer lies not in a machine, but in a method: a set of methods accumulated over thousands of years specifically to recognize the past in fragments. Learn them, and you will not only understand history books better. You will also see through much of the ``it's said that'' on your phone.

\section{``Dragon Bones'' in a Pharmacy}
Let us enter a scene. The scene itself is a legendary opening to the question of how we know history.

Time: autumn 1899. Later accounts give several versions of the details; what follows is the most widely circulated. Place: a traditional medicine shop in Beijing. The air is thick with medicinal smells, the sweetness of licorice, the bitterness of dried tangerine peel, the sharp smoke of mugwort, all mingled together. Behind the counter, an assistant swings a small copper hammer, breaking an ingredient into pieces before weighing it.

The ingredient is called ``dragon bone,'' said to be fossilized bones of large ancient animals and used in Chinese medicine to calm the mind. The customer is Wang Yirong, head of the imperial academy, roughly the president of the country's highest educational institution in today's terms. He is also an epigrapher who has spent his life studying inscriptions on ancient bronzes.

As the assistant breaks open a piece of ``dragon bone,'' Wang Yirong's gaze stops.

There are marks on the fragment. Not insect damage, not cracks, but \textbf{writing}. Cut with a blade, the strokes thin and firm, arranged in rows. After a lifetime studying bronze inscriptions, he recognizes at once that their style is older than any he knows. What the assistant sees as worthless scraps of medicine suddenly becomes, to Wang, a library.

Events then snowball. Wang begins paying high prices for inscribed dragon bones. Tracing their origins reveals that most come from around Xiaotun village in Anyang, Henan, where farmers turn them up while plowing and sell them as medicinal ingredients. Wang dies soon afterward amid warfare, and others take up the research. In 1928, archaeologists begin formal excavations at Xiaotun: the world-famous excavation of Yinxu, the late Shang capital. Those ``dragon bones'' are oracle bones used by Shang kings more than three thousand years earlier. Before war, ask a question; before the harvest, ask another; whether the queen's child will be a boy or a girl, ask that too. Questions are carved on turtle shells and animal bones, heat produces cracks, and the cracks are read for good or ill fortune.

Something almost unbelievable happens. The Shang dynasty, previously present only in ancient books and whose very existence some had doubted, suddenly reaches out of the soil with tens of thousands of inscribed bones and begins to speak.

Remember this medicine shop. Every question in this chapter is hidden in those two seconds when someone with a trained eye stares at an unwanted scrap of bone and recognizes writing. \textbf{Seeing evidence takes skill.} In the assistant's hands the same thing is medicine; in Wang Yirong's it is history. The difference lies not in the object, but in the eyes.

\section{What Makes the Past Count?}
Let us lay out the conflict without rushing to an answer.

Open any history textbook and you meet a remarkably confident style: ``In 1046 BCE, King Wu attacked King Zhou.'' ``Pompeii was destroyed in 79 CE.'' ``Tang Chang'an had more than a million inhabitants.'' The sentences are brisk and clear, like a reporter's dispatch written on the day.

Think for a moment, though, and something feels wrong. No reporter stood by when King Wu attacked. No camera watched Pompeii disappear beneath volcanic ash. Chang'an's census sheets rotted away long ago. None of these statements comes from our own eyewitness experience. All were \textbf{pieced together by later people working with incomplete evidence}.

Two camps begin to argue, and have kept arguing for more than two thousand years.

One says that something pieced together cannot be trusted. The evidence is fragmentary, its writers had their own interests, and thousands of years of war, fire, flood, and forgetting stand between us and them. Who knows which fragment has been misplaced? Who even knows whether it is a genuine fragment rather than a later forgery? This camp's sharpest claim is that history is merely ``lies everyone agrees to believe.''

The other says: What is wrong with fragments? Criminal investigations use them too. With sound methods and enough evidence from varied sources that check against one another, the assembled picture is reliable. Before oracle bones were discovered, some doubted that the Shang dynasty was real. Afterward, the kings' names matched one by one, leaving the doubters speechless. This camp believes that although we cannot see the past directly, we can \textbf{get closer to it}.

The real difficulty lies between these partly justified positions: \textbf{With no recording to watch, yet no license to invent, what entitles us to say that something ``happened in history''?}

The question is not abstract at all. It determines which textbook statements you should trust and which short-video ``little-known facts'' you should not. It affects whether a decades-old court case can be reopened, and how much weight ``Your grandfather says that back then...'' should carry at a family gathering.

Before reading on, take a position. Do you think we can truly ``know'' a past that has no recordings?

\section{Hear Out Both Kinds of Doubt}
First, let us give the skeptics their full case. They are often dismissed as argumentative troublemakers, and that is unfair.

Their reasons grow stronger layer by layer. First, \textbf{sources are written by living people, and living people have positions}. Victors can blacken the defeated while writing history; rulers can omit shameful acts while recording achievements. Xiang Yu lost the Chu--Han struggle, but who wrote the ``Basic Annals of Xiang Yu'' we read today? Sima Qian, a Han court historian, an employee of the winning side. Although his account is fairly generous, the very fact that the loser's story is told by the victor's historian deserves scrutiny. Second, \textbf{sources can die}. Wars burn books, floods soak archives, insects eat pages, and successive courts deliberately destroy ``unsuitable'' records. The ancient documents we can read are survivors of repeated selection and countless disasters. In statistical language, they form a severely biased sample. Third, \textbf{most living people were never recorded at all}. Few people in antiquity could write. A Tang farmer's whole life, what he grew, what he feared, whom he loved, how he died, never caused a document's writer to pause over him. The names crowding history books belong to emperors, generals, and ministers. The silent majority left not even names. Equating history with the history of those recorded is like drawing an entire nautical chart from the positions of a few lighthouses.

Are these strong reasons? Very strong. So strong that any honest historian must accept most of them.

But hear out the other side too, with its case stated just as fully.

\textbf{First, a forgery may fool one kind of evidence, but not every kind.} You can invent a text, but it is hard to invent matching sites, artifacts, coins, place names, dialect loanwords, and pollen in soil layers all at once. The more kinds of evidence, the greater the cost of lying. This is the corroboration we will discuss in detail. \textbf{Second, a writer's position does not mean every sentence is a lie.} A position shapes selection and wording, but much hard information, dates, offices, prices, population figures, offers little incentive to lie and many ways to check. \textbf{Third, silence itself can leave a shape.} Ordinary people left no writing, but they left village sites, graves, charcoal in hearths, and signs of strain on bones. Archaeologists cannot read their words, but can read their life spans, nutrition, diseases, and labor. Silence is not emptiness. It is another kind of evidence, just harder to read.

Add a footnote to these two voices from another part of the world. On ancient Egyptian temple walls, pharaohs carved spectacular accounts of victories and simply omitted defeats. The voice of the center was loud and bright. Yet in workers' villages excavated near the desert's edge, remains in bread molds and absence records on inscribed stone flakes, with reasons including ``burying his brother,'' tell quite another story. The rulers' and victors' monuments are real; the crumbs left by marginal people are real too. The historian's skill is to make these voices check one another, not to hear only the louder one.

So the real dispute is not whether to believe history, but \textbf{how much credit each kind of evidence deserves}. That is exactly what the next section's tool helps you judge.

\section{Thinking Bridge: Give Evidence a Checkup}
A doctor does not simply declare someone ``healthy'' or ``ill.'' The doctor checks one thing after another. Historians treat evidence similarly. This checkup has six questions, and you can take them unchanged to any information claiming to come from the past.

\textbf{Question zero: How many steps from the event is it?}

A primary source is a participant's account or a contemporary object: an inscription made on site, a letter written at the time, a participant's diary, an artifact left by its user. Secondary sources are later compilations and studies: historical works, textbooks, documentaries, and this very book.

Now immediately discard two widespread misconceptions.

Misconception one: primary source equals absolutely reliable. Wrong. Here is a fine counterexample. When Vesuvius erupted in 79 CE, the Roman Pliny the Younger was across the bay. Later he wrote two letters to the historian Tacitus describing that day and his uncle's death. These are genuine primary sources. Read closely, however, and the uncle in the letters is fearless and observant, while ``I'' is composed, studious, and unshaken by crisis. The letters were very likely written with posterity in mind; image management runs between their lines. Eyewitness memory can fail, pride can inflate it, and circumstances can encourage embellishment. A primary source is close to the scene, not necessarily close to the truth.

Misconception two: secondary research equals hearsay. Wrong. When Sima Qian wrote about the Five Emperors and the Xia and Shang dynasties, the events were already more than a thousand years old. His work was thoroughly secondary. Yet he visited sites widely, checked documents, compared traditions, and honestly recorded several versions where matters remained unclear. A rigorous secondary study often has more value than a casually written primary record. The number of steps tells us about distance, not directly about reliability. The following five questions test reliability.

\textbf{Question one: Who left it?} (Author.) The oracle-bone inscriptions were made by the Shang king's diviners; a letter home was written by a homesick soldier. The author's identity shapes what they could and could not see.

\textbf{Question two: Who was it for?} (Audience.) An imperial edict addresses the realm, a diary the writer, a eulogy the mourners. The same person speaks differently to different audiences. That Pliny's letters addressed future readers already explains a great deal.

\textbf{Question three: Why was it made?} (Purpose.) To keep a record? To spread propaganda? To claim credit? To avoid tax? The more directly useful the purpose, the greater the temptation to ``process'' the facts.

\textbf{Question four: When was it made?} (Date.) A record made during an event and a memoir written fifty years later must receive different treatment. Memory quietly rewrites history all by itself.

\textbf{Question five: How did it survive until today?} (Preservation.) Every stop between an object's creation and your encounter with it can introduce trouble. Copyists make mistakes, translations distort, anthologies omit, excavated objects break. Online information is no different: how many hands a screenshot passed through and what was cropped away form its preservation history.

The last question is the easiest to overlook and often the most consequential. An ancient book survives two thousand years not through its original manuscript, long since gone, but through generations of handwritten copies. Each copying creates another chance for error: wrong characters, omitted lines, marginal notes copied into the body. All have happened. The ``original text'' you read today is a surviving version after countless copyings. Scholars have developed elaborate ways to compare versions, recording which character appears in which form in which copy. It sounds dull, but it is the firmest foundation beneath the question of why we should trust words from two thousand years ago.

Now give your grandfather's half coin a complete checkup. It is a primary source, a contemporary object. Its maker was an unnamed mint worker. Its audience was everyone who used it. Its purpose was circulation, not recordkeeping, making it reliable on contemporary prices and casting techniques but silent about who spent it. Its date must be checked against its inscription and form. Its preservation history is casting $\rightarrow$ circulation $\rightarrow$ loss $\rightarrow$ burial $\rightarrow$ discovery $\rightarrow$ a drawer. Once examined, the silent coin has already ``confessed'' quite a lot, while marking what it can never reveal. Knowing what evidence \textbf{cannot} prove is as important as knowing what it can.

\section{Even Confucius Lacked Sources}
Now read a short primary source. You will discover that someone was already writing about the frustration of insufficient evidence two and a half millennia ago.

In the ``Eight Rows'' chapter of the Analects, Confucius says:

\begin{quote}
The Master said, ``I can speak of the rites of Xia, but Qi cannot sufficiently attest to them. I can speak of the rites of Yin, but Song cannot sufficiently attest to them. This is because written records and knowledgeable people are insufficient. Were they sufficient, I could substantiate them.''
\end{quote}
Roughly: I can discuss the Xia dynasty's ritual system, but the state of Qi, home to the Xia people's descendants, lacks enough evidence to verify it. I can discuss the Shang ritual system, but Song, home to the Shang descendants, also lacks enough evidence. Too few written sources and people familiar with the old ways survive. If there were enough, I could verify it.

Weigh these words. Discussing ritual, Confucius does not show off his learning. He first states the evidence's condition: what he can describe, what he \textbf{cannot verify}, and why verification fails, namely insufficient records and knowledgeable people. He does not turn tradition into fact or force an incomplete story into completeness. In one statement, knowing and knowing insecurely are clearly separated. That is awareness of evidence. Notice, too, how the term wenxian is divided into two kinds of source: wen, written materials, and xian, elders familiar with old affairs. Writing plus oral testimony: a checklist from 2,500 years ago is not fundamentally different from a historian's today.

Confucius lacked sources. We are somewhat more fortunate. Historians today have six ``families of evidence.''

\textbf{First, writing.} Oracle bones, inscriptions, bamboo and wooden slips, paper books. The most dramatic example comes from Egypt. The Rosetta Stone carries the same decree in three scripts: two ancient Egyptian scripts nobody could still read, alongside readable ancient Greek. This comparative key helped scholars open the door to ancient Egyptian civilization. Across the world in Central America, Maya people carved kings' names and dates of campaigns on stone steps and monuments. After their writing had been lost for centuries and was deciphered again, forest dwellers once dismissed as ``primitive'' suddenly possessed a whole dynastic history written by themselves.

\textbf{Second, artifacts.} Potsherds, coins, tools, ornaments. In 1991, hikers discovered Ötzi the Iceman in an Alpine glacier, a body more than five thousand years old accompanied by a copper axe, a bow and arrows, and tinder fungus. One person's equipment forms a miniature exhibition of an age.

\textbf{Third, sites.} Cities, houses, graves, and rubbish pits. Do not underestimate rubbish pits. Discarded pottery and bones can be particularly honest because nobody cherished or manipulated them.

\textbf{Fourth, images.} Rock paintings, murals, carved pictorial stones. Han pictorial stones show cooking, horses and carriages, and feasts: ``photographs'' left by kitchens and stables that never wrote words.

\textbf{Fifth, oral testimony.} West Africa has professional keepers of memory called griots. They can sing for days and nights about Sundiata, the founder of the Mali Empire, preserving centuries of royal genealogy through successive generations of memorization. Oral accounts change, but their skeletons, genealogies, place names, major events, can be remarkably stable. Something similar happens on the Eurasian grasslands. Tibetan storytellers pass on the Epic of King Gesar orally, a work longer than any fixed written epic and still growing. How should oral sources be used? Historians treat them as they treat writing: neither discard them as ``mere legends'' nor believe everything because it has been handed down. They compare them with sites, artifacts, and other peoples' records, retaining what matches and reserving judgment on what does not.

\textbf{Sixth, language and genes.} Language keeps accounts. Chinese words such as putao, grape, and shizi, lion, came from the Western Regions. The words themselves are receipts from ancient trade routes. Genes keep accounts too. Europeans and Asians today still carry small amounts of Neanderthal ancestry. Encounters tens of thousands of years ago were written into bodies without the people involved ever knowing.

Evidence never appears alone. Different sources testify for one another and undermine one another. The next section puts that point to work.

\section{How Bad Was King Zhou?}
Now consider how many explanations can grow from the same ``fact.''

In traditional histories, the last Shang ruler, King Zhou, is a textbook tyrant: pools of wine, forests of meat, torture on a burning bronze pillar, the cutting out of loyal minister Bi Gan's heart. King Wu of Zhou's campaign against him therefore becomes an act of heavenly justice. These accounts mainly come from documents written in or after the Zhou dynasty. Notice who wrote them: the side that destroyed the Shang and those who inherited its perspective.

Yet Confucius's disciple Zigong had doubts. In the ``Zizhang'' chapter of the Analects he says:

\begin{quote}
Zhou's wickedness was not so extreme as this. Thus the gentleman dislikes dwelling downstream, where all the world's evil gathers upon him.
\end{quote}
Roughly: King Zhou was bad, but not as bad as the stories claim. Once someone is branded a bad person, every evil in the world is piled onto them. Zigong identifies a rule still operating today: reputations snowball, and bad ones especially fast.

In the twentieth century, the historian Gu Jiegang turned this suspicion into a systematic account. Broadly, ancient traditions accumulate in layers. Later ages add more deeds to ancient figures, making their stories ever fuller: the famous theory of layered accretion. To demonstrate it, he did painstaking work. He copied each accusation against King Zhou from successive historical texts into a chronological table, marking the earliest book in which each appeared. The result was striking. Early Western Zhou documents offered only a few, fairly general accusations. Only in books written centuries later did the vivid scenes appear: wine pools, meat forests, the burning pillar, the pit of venomous creatures. The later the period, the more accusations and the more vivid the details. Details of a real event should grow hazier with time. Their growing clarity instead points to \textbf{reworking}.

At least three explanations of how bad King Zhou really was now stand before us, and they are not equivalent.

\textbf{Explanation one: He really was cruel.} The Zhou attacked Shang under the banner of relieving the people and punishing wrongdoing, then repeatedly proclaimed the former ruler's evils. This helped legitimize the new regime, but effective propaganda generally needs some factual foundation. This explanation relies on ``no smoke without fire.'' Its limit is that it cannot explain why King Zhou's crimes multiply in later ages. A person cannot keep committing new crimes centuries after death.

\textbf{Explanation two: Stories accumulated in layers.} The real King Zhou may have been merely a forceful, unsuccessful last ruler. Each later generation of storytellers and each political commentator needing a negative example added another stroke, finally producing the supreme tyrant. This elegantly explains the strange multiplication of crimes over time. Its limit is that it cannot establish the nature of the earliest underlying layer, and can easily slide too far into ``all ancient history is invented.''

\textbf{Explanation three: Each accounts for half.} King Zhou's rule did contain cruelty, the ``fire'' recognized by the first explanation, while particular stories underwent repeated reworking, the ``smoke'' recognized by the second. This fits the existing evidence best. Its limit is that ``half and half'' is a vague verdict. Every detail must still be investigated separately; there is no once-and-for-all formula.

How do we weigh these explanations? Through corroboration. In 1917, the scholar Wang Guowei did something groundbreaking: he compared newly excavated oracle-bone inscriptions with the list of Shang kings in the ``Basic Annals of Yin'' in the Records of the Grand Historian. Generation by generation, name by name, they matched. The inscriptions even corrected some misplaced generations in the written account. This is the famous ``dual-evidence method'': material on paper and material underground checking one another.

That corroboration adds substantial weight to the genealogy in the Records. It is no longer merely reported; it has supporting evidence. But notice the weight's limits. Oracle inscriptions can confirm the list of Shang kings, not the details of wine pools and meat forests. A reliable genealogy does not make every story reliable. \textbf{Corroboration raises confidence in the chain of evidence; it does not exempt every link from inspection.} That sense of proportion is what ``weight of evidence'' means.

Consider a counterexample that is easy to misjudge, showing how the same evidence can overturn an interpretation. When Ötzi was first found, most experts thought he was a shepherd who had lost his way and frozen in a blizzard, a plausible account. Then, in 2001, a CT scan revealed an arrowhead in his left shoulder, rewriting the cause of death: someone had shot him from behind. The same body, the same primary source, but new evidence overturned the explanation. That does not show history to be unreliable. On the contrary, this is how it earns reliability: explanations remain open to new evidence.

\section{Further Challenge: Between Fact and Story}

We have reached this chapter's weightiest distinction: \textbf{historical fact} and \textbf{historical narrative} are not the same thing.

The encirclement at Gaixia, Xiang Yu's defeat, his suicide at the Wu River: these belong to the factual level, a broadly undisputed framework supported by the Records and later documents. But close your eyes and your Xiang Yu is probably not a defeated warlord. He is a tragic hero, strong enough to lift a great cauldron, defeated by fortune, parting from Lady Yu and his horse Wuzhui, refusing to return east across the river. This Xiang Yu belongs to the narrative level, a building Sima Qian's pen raised upon historical facts. Tell the same framework differently and you can build something else: a headstrong warlord who massacred cities, buried surrendered soldiers alive, and was finally brought to account by his age. Same framework, different building.

The historian Hayden White studied this specifically. Broadly, he argued that historical writing does not simply state the past. It inevitably \textbf{makes the past into a kind of story}, tragedy, comedy, epic, or satire, borrowing its plot patterns from literature. This provocative view need not be accepted wholesale to acknowledge the crucial point it touches: selecting facts, choosing their order, deciding where to end are themselves acts of interpretation, not neutral transportation.

For example, the four years of the Chu--Han struggle are a tragedy from Xiang Yu's viewpoint, a founding epic from Liu Bang's, and an utter catastrophe from the perspective of a Pei County farmer conscripted three times who never returned home. \textbf{Three buildings made from the same bricks.} Historical fact answers what happened. Historical narrative answers what it meant, and that answer always includes the narrator's position, genre, and intended readers.

You can test the distinction nearby. When the same person dies, a eulogy presents a kind elder, a court judgment, if an inheritance dispute arises, a party to rights and obligations, and a genealogy a member of a particular generation and family branch. All three contain facts, yet construct different buildings. Think also of news. Official reports from opposing sides of one conflict may read like accounts of two wars, even when most facts cited by both are true. Narrative's magic lies not necessarily in lying, but in \textbf{selection and arrangement}: which bricks were chosen, which laid first, and where the last one ends.

Now place silent evidence within this framework and a deeper layer appears. Silence is not merely a missing piece. It is \textbf{structural}. Few could write, so few left narratives. Those in power decided what deserved recording, so records largely concerned power's interests. Farmers, women, enslaved people, and the defeated formed the great majority in fact, yet were nearly invisible in narrative. This is not one historian's malice, but an imbalance in the whole recording mechanism. Once you see it, any historical passage prompts an additional question: Who is telling this to whom? Where are those who never appear?

Distinguishing fact from narrative does not teach nihilism. It does not mean, ``It's all storytelling anyway, so don't take it seriously.'' Quite the opposite: it gives you two rulers. One measures facts: Is there evidence? How close to the event? Does it pass its checkup? The other measures narrative: Who is speaking? What kind of drama have they made? Who is absent? Only with both are you a complete reader of history.

\section{Evidence in the Screenshot Age}
This question, more than two thousand years old, goes to court on your phone every day.

First consider the historical side. Short-video platforms produce astonishing amounts of ``little-known history'': ``The First Emperor was actually...'' ``What archaeologists dare not reveal...'' Apply our six-question checkup and most collapse immediately. Author? A marketing account. Audience? Online traffic. Purpose? More followers. Date? Invented yesterday. Preservation? Copied from forum to forum, passed through eighteen hands. Ancient sources may be scarce, but today's fake sources are manufactured by the ton.

Now consider the present. We once said, ``Pictures prove it.'' Today images are among the most questionable evidence. Screenshots can be forged, faces replaced in videos, and footage of an elder tearfully remembering the past may have both a generated face and a generated voice. This forces us into a position resembling that of earlier ages: evidence needs examination, not merely eyes. An online screenshot's preservation history is its forwarding chain. How many hands? Has context been cropped? Does the original post survive? When judging a screenshot, you do the same work as Sima Qian comparing traditions or Wang Guowei checking oracle inscriptions against the Records. Only the tools have changed, from knife and brush to a search box.

The method is useful at school every day too. A class chat says, ``The monthly exam is next week,'' with ``the other class said so'' as its source. Start with question zero: how many steps from the event? A school announcement is primary; a retelling is secondary; ``I heard'' may not even qualify as secondary. At a family gathering, elders argue heatedly over whether great-grandfather traded goods or farmed. This is oral history in action. Everyone's memory is genuine, and everyone's memory has been quietly rewritten over decades. Even the ``story of three generations of my family'' in your essay is, strictly speaking, an unexamined primary source. Not one of the historian's six questions is wasted in such settings.

There is a more personal layer: \textbf{you are making historical sources yourself}. Chats, social posts, shopping lists, and location trails are being preserved at a density ancient people could not imagine. A Tang farmer's entire trace in history might be one bone in a grave. You create enough records in a day for a future person to write your biography, if they can access them and read them correctly. The old problem was the silent majority. The future's may be the reverse: so much survives that nobody can read it all. Storage formats expire, and two thousand years on a hard drive may be less survivable than time on an oracle bone.

So this chapter's method was never merely a historian's craft. Examine each piece of information, distinguish fact from narrative, and ask where the silent people are. These are everyday self-defense skills for everyone in the information age.

\section{Over to You: A Person Who Left No Record}
Return to the half coin in your grandfather's drawer.

It cannot prove who spent it. The person who lost it, perhaps a peddler heading to market, perhaps a child who had saved for half a year, has no name, no biography, no line of writing that remembers them. Now answer this chapter's final question, and you must give reasons:

\textbf{Does someone who left no record throughout their life have a ``history''?}

Before answering, weigh what is in your hand once more.

Those who say no have strong reasons. History speaks through evidence. About someone who left no trace, we know nothing, and ignorance is not history. Even Confucius acknowledged that insufficient sources mean insufficient corroboration. Admitting we do not know is more honest than pretending.

Those who say yes have strong reasons too. The fields they worked, houses they built, coins they spent, and genes they passed on are real traces. Their silence was not their choice, but the recording system's choice. Turning ``not recorded'' into ``without history'' makes them pay a second time for that system's absence.

A harder question waits. If future archaeologists one day dig up your phone, its chips intact and data readable, could they reconstruct the ``real you''? Or would they, like us today, hold a pile of evidence and write a dozen versions of you, each both right and wrong?

There is no standard answer. But notice: from the moment you seriously consider this question, you are no longer only a reader of history. You have begun to think like a historian.
